\chapter{Conclusion}
\label{ch:conclusion}

This project built a teleoperation system for two seven-degree-of-freedom arms
mounted on a worn backpack frame, commanded from an instrumented mannequin
master, with a safety architecture appropriate to a manipulator standing
beside a person and a shared-autonomy layer above it.

\section{Against the objectives}

\emph{A master, built and characterised.} Fourteen potentiometer channels, two
inertial sensors, two force sensors and two buttons on a single
microcontroller, emitting a text frame at \SI{100}{\hertz}. Characterisation
found six of fourteen channels usable and reversed two long-standing verdicts,
both of which had come from computing statistics across transmitted rows
rather than distinct sensor updates.

\emph{A mapping that degrades predictably.} Position in spherical form with
each component taken from the sensor that observes it best. Vertical and
fore-and-aft motion track correctly on both arms. Lateral does not, and the
measurements show why no axis map or recalibration can fix it. The cost of
each channel reduction is quantified on \num{20430} recorded frames, and the
endpoint of the ladder reports position unavailable rather than inventing one.

\emph{A safety architecture validated by injection.} Fourteen of fourteen
injected faults handled. The two that took several attempts both failed for
the same underlying reason, which is that a mechanism was keyed on the wrong
observable: a dead-man on message arrival cannot detect a source that keeps
publishing stale data.

\emph{A workspace characterisation.} The two arms share no reachable
workspace under the as-built mount, which blocks inter-arm handover for
geometric reasons no software change can address. Asked four separate ways,
the working space in front of the wearer's own chest is unavailable at any
surface height or distance, and deleting the wearer's arms from the model
moves the boundary by at most \SI{75}{\milli\metre}, which identifies the
torso rather than the limbs as the constraint. The quantified mount
modification was subsequently applied as one deliberate pass, raising the
clearance ceiling from \SI{0.1610}{\metre} to \SI{0.2202}{\metre} and bringing
the task's placement targets from \SI{0.530}{\metre} to \SI{0.290}{\metre} off
the centre line.

\emph{A verification methodology.} Ten repeats over densified paths rather
than single calls at endpoints, and an instrument validated against a known
answer before a surprising measurement taken with it is reported. The
five-task verification reproduced exactly across runs, which is the first
reachability figure in this project to do so.

\emph{Motion generation and transfer.} The per-joint step limiter that
constituted the entire motion generation was replaced with a jerk-limited,
phase-synchronised generator, removing \SI{51.3}{\milli\metre} of unrequested
Cartesian excursion on a platform whose capture gate is
\SI{30}{\milli\metre}, and bringing peak commanded joint speed within the
manufacturer's declared limit for the first time. The simulation-to-hardware
gap was measured rather than assumed: every joint parks \SI{0.305}{\degree}
short on the side it approached from, which one scalar per arm removes three
quarters of.

\emph{Measurement in place of declaration.} A calibration stage measures the
working surface and the objects on it, recovering six objects with every
centre within \SI{1}{\milli\metre}, and hands the result to a planner that had
previously known about nothing except the wearer.

\emph{Identification of what blocks progress.} Seven items, listed in
\cref{ch:discussion} in the order they should be addressed. The first is two
soldered joints.

\section{What was not achieved}

No result in this thesis was validated on a physical Kinova arm carrying a
load. One session drove a single arm from its own camera and produced three
findings, all of them about calibration and none of them a performance result
(\cref{sec:hardware}). The task set is verified geometrically and has never
been performed with a wearer in the harness. The two-person study
the architecture was designed around was not run, and could not have been:
the direct teleoperation baseline depends on channels that are currently
incoherent.

The planning report's haptic objective was not delivered and, as
\cref{ch:discussion} argues, is not deliverable on this platform through the
available control interface.

\section{Future work, in the order it should be done}

\begin{enumerate}
  \item \textbf{Repair left j2 and j4, and re-run the channel baseline after
    each attempt.} These two are named by the observability regression rather
    than chosen: with both frozen the left arm retains no reach observable at
    all, and repairing any other channel will not restore one. Re-running
    after each attempt rather than once at the end catches a repair that
    breaks a working channel.
  \item \textbf{Correct the filter seed in the firmware} so the potentiometers
    do not spend \SI{440}{\milli\second} settling from a railed value after
    every board reset. The change is one line and does not alter the filter's
    steady state.
  \item \textbf{Read the right arm's home joint angles from hardware.} They
    have never been read, and the workspace geometry depends on them.
  \item \textbf{Bring up a second inertial sensor per arm on the secondary
    bus.} The elbow angle then comes from the relative orientation of two
    segments with no potentiometer, and the shared yaw ambiguity cancels in a
    relative measurement, which is why two sensors can replace the elbow and
    one cannot. It costs two potentiometer channels and reaches
    \SI{63.8}{\milli\metre} on a single remaining pot.
  \item \textbf{Calibrate the wrist camera against the arm.} It is the largest
    term in the positioning budget, it is wrong by roughly \SI{8.5}{\degree}
    in a way that rotates with the wrist, and it is the reason no open-loop
    grasp has landed inside the capture gate. It needs the hardware and a
    target, and it is the cheapest high-value measurement remaining.
  \item \textbf{Calibrate the scene colour camera}, which has never been
    calibrated in this project and whose eight refusing stages
    (\cref{ch:vision}) all refuse for the same missing intrinsics. It is a
    half-hour with a printed board and it converts a camera that can only
    report pixels into one that can report metres.
  \item \textbf{Re-run the stage-by-stage vision session with the arms
    actually at the pick pose.} The session of \cref{ch:vision} measured the
    arms at \SI{10.67}{\degree} and \SI{13.53}{\degree} from the reference
    configuration, so none of its figures may be captioned as a picture of
    the pick pose, and the numbers describe a viewpoint near it rather than
    the one the task uses. Nothing needs to be built; the arms need to be
    parked and the same command re-run.
  \item \textbf{Make the clearance check measure link geometry.} At present it
    samples link origins and returns a comfortable number for a pose in which
    the arm is touching the mannequin. This is a safety defect rather than an
    incompleteness, and it should be closed before anybody wears the rig.
  \item \textbf{Put the perception latency into the clearance floor.} One
    parameter, and it converts a margin that means little into one that means
    what it claims.
  \item \textbf{Re-park the right arm and re-derive the mount again if a
    shared workspace is wanted.} The applied mount change improved clearance
    and brought the task inboard; it did not make the two reachable sets meet,
    and nothing in software will. It should again be one deliberate pass,
    because the home pose, the anchor and every clearance figure move
    together.
  \item \textbf{Validate the safety layer on hardware}, starting with the
    session recovery path, which is the riskiest because the arm permits
    exactly one session.
  \item \textbf{Run the two-person study.} The exchange rate between operator
    performance and wearer trust is the open question this architecture was
    built to ask, and it remains open.
\end{enumerate}

\section{Closing}

The most transferable outcome of this work is not the mapping or the safety
architecture. It is the discipline that produced them. Twenty-five times a
result that looked like a system failure turned out to be a defect in the
measuring tool, and fourteen of those were caught only because something else
already had the answer. On a robot that stands beside a person, the
difference between a measurement and an artefact is not an academic
distinction, and the habit of establishing which one is in hand before acting
on it is the part of this project most worth carrying forward.
